% Trabalho academico ABNT (abnTeX2). Marcadores <<CHAVE>> sao trocados por doclib.fill_template();
% o corpo do texto vai em <<TEXTO>> (LaTeX: \chapter{...}, \section{...}, \citeonline{...}, \cite{...}).
% Compile com doclib.latex_abnt("trabalho.tex") ou: latexmk -r /app/templates/abnt/latexmkrc trabalho.tex
\documentclass[12pt, oneside, a4paper, chapter=TITLE, section=TITLE, sumario=abnt-6027-2012, english, brazil]{abntex2}

\usepackage{lmodern}
\usepackage[T1]{fontenc}
\usepackage[utf8]{inputenc}
\usepackage{indentfirst}
\usepackage{color}
\usepackage{graphicx}
\usepackage{microtype}
\usepackage{booktabs}
\usepackage[alf]{abntex2cite}

% Dados da capa e da folha de rosto
\titulo{<<TITULO>>}
\autor{<<AUTOR>>}
\local{<<LOCAL>>}
\data{<<ANO>>}
\orientador{<<ORIENTADOR>>}
\instituicao{<<INSTITUICAO>>}
\tipotrabalho{<<TIPO_TRABALHO>>}
\preambulo{<<PREAMBULO>>}

\setlength{\parindent}{1.25cm}
\setlength{\parskip}{0pt}
\SingleSpacing
\OnehalfSpacing
\makeindex

\begin{document}
\selectlanguage{brazil}
\frenchspacing

% ---------- Elementos pre-textuais ----------
\imprimircapa
\imprimirfolhaderosto*

\begin{resumo}
<<RESUMO>>

\textbf{Palavras-chave}: <<PALAVRAS_CHAVE>>.
\end{resumo}

\begin{resumo}[Abstract]
\begin{otherlanguage*}{english}
<<ABSTRACT>>

\vspace{\onelineskip}
\noindent\textbf{Keywords}: <<KEYWORDS>>.
\end{otherlanguage*}
\end{resumo}

\pdfbookmark[0]{\contentsname}{toc}
\tableofcontents*
\cleardoublepage

% ---------- Elementos textuais ----------
\textual

<<TEXTO>>

% ---------- Elementos pos-textuais ----------
\postextual
\bibliography{referencias}

\end{document}
